% Level 1, Week 6 -- homework
\homeworkweek{6}{Review Studio and Checkpoint}{15--20 minutes. Do Core first; Challenge and Creative are optional.}
\namedateline

\begin{core}[Remix Studio (about 10 minutes)]
\textbf{Part 1: Blocks to words (paper, 4 minutes).} Use your card sort to write the Python line for each block. Guess if you are not sure: guessing is good practice!

\begin{center}
\renewcommand{\arraystretch}{1.4}
\begin{tabular}{@{}m{6.2cm}m{8cm}@{}}
\begin{scratch}[scale=0.75]\blockmove{move \ovalnum{50} steps}\end{scratch} & \rule{7cm}{0.4pt} \\
\begin{scratch}[scale=0.75]\blockmove{turn \turnleft{} \ovalnum{90} degrees}\end{scratch} & \rule{7cm}{0.4pt} \\
\begin{scratch}[scale=0.75]\blocklook{say \ovalnum{Hi!}}\end{scratch} & \rule{7cm}{0.4pt} \\
\begin{scratch}[scale=0.75]\blockrepeat{repeat \ovalnum{3}}{\blockspace}\end{scratch} & \rule{7cm}{0.4pt} \\
\begin{scratch}[scale=0.75]\blockvariable{set \selectmenu{lives} to \ovalnum{3}}\end{scratch} & \rule{7cm}{0.4pt} \\
\end{tabular}
\end{center}

\textbf{Part 2: Level up a project (Scratch, 6 minutes).} Open one of your projects from weeks 1--5 (\menu{File} $\rightarrow$ \menu{Load from your computer}). Add \textbf{one} improvement from this menu, test it, and save it with a new name ending in \texttt{-v2}:
a sound effect \;\textbullet\; a new backdrop \;\textbullet\; a score or lives variable \;\textbullet\; a key that does something \;\textbullet\; a My Block \;\textbullet\; a message between two sprites.

My project: \rule{4cm}{0.4pt}\quad What I added: \rule{6cm}{0.4pt}
\end{core}

\begin{challenge}[A bug for a friend]
Build a tiny project (8 blocks or fewer) with \textbf{one} bug hidden inside. On a card, write what it \textbf{should} do, and draw or list its blocks. Bring the card to the next session: a classmate will try to find your bug!
\par\textbf{It should:} \answerlines{1}
\textbf{The bug (keep this secret, fold it over):} \answerlines{1}
\end{challenge}

\begin{creative}[Remix with a new theme]
Make a remix of Apple Catch with a new theme, for example catching stars in space, or fish in the sea. Change the sprites, the backdrop and the sounds (from the Scratch library or drawn by you). Add a credits comment: \enquote{Remix of Apple Catch by} + your nickname, + \enquote{sprites from the Scratch library}.
\end{creative}

\begin{selfcheck}
  \item I wrote a Python line for every block in Part 1.
  \item I added one improvement to an old project and it works.
  \item I saved the new version with \texttt{-v2} at the end of the name.
  \item My remix (if I made one) has a credits comment.
\end{selfcheck}
